The authors of \cite{cigarran_twitter_fca_2016} apply \ac{fca} to a Twitter dataset to extract insights into the topics present.
% different: Use terms as attributes -> like frequency based and not semantic based (like our topic model)
In their methodology, tweets serve as objects, while words occurring within the tweets act as attributes.
They assume that tweets forming a formal concept, i.e., those sharing the same terminology, correspond to the same topic.
Upper concepts in the resulting concept lattice represent general topics,
whereas lower-level concepts reflect more specific themes.
% different: Discard not only bottom concepts and do not use TITANIC
However, a key limitation of their approach is the excessive number of concepts generated, making interpretation challenging.
To mitigate this, they introduce upper and lower frequency thresholds,
discarding terms whose occurrence falls outside these predefined bounds,
arguing that the former are too general and the latter too specific for meaningful analysis.
While their approach provides valuable means as how to derive topic structures from text data,
it relies purely on term frequency as attributes rather than
leveraging the semantic relationships inherent in topic models, as we do in this work.

The authors of \cite{wang_topic_feature_2021} employ \ac{fca} to structure topic representations derived from a topic model.
They construct two matrices:
One capturing document-topic probabilities obtained through their ILDA model and
another representing topic-word probabilities, restricted to the top ten words per topic.
% different number of topic words, but same idea
% Their approach considers only the top ten words per topic,
% thereby representing topics in a manner comparable to our work.
% different: Manual topic categorization
Unlike our approach, which automatically derives hierarchical relationships,
their method relies on manual categorization of topics into broader thematic groups.
These manually curated categories are subsequently visualized using a \ac{fca}-inspired order diagrams to reflect topic generalization.
In contrast, our approach seeks to infer hierarchical topic structures automatically, without requiring manual intervention.

The authors of \cite{wang_topic_hierarchy_2023} also leverage \ac{fca} to structure topic hierarchies.
% different Topic Model
They utilize LDA for topic modeling, encoding documents based on the distribution of topics.
% same: use threshold
Similar to our approach, they introduce a thresholding mechanism to retain only relevant topics per document.
However, a major divergence lies in their reliance on HowNet, a linguistic knowledge base.
% to extract for each word senses and a set of words describing the senses.
% This enables them to construct a document–sememes concept lattice,
% which supports fine-grained document retrieval.
% Documents are considered to be objects whereas sememes are treated as attributes.
% different goal
Their primary objective is to facilitate precise querying,
using a concept lattice to provide a more nuanced representation of documents beyond mere word-level indexing.
% different rules to reduce size
To manage the complexity of the lattice, they eliminate concepts whose attribute sets exhibit excessively high or low support.
While effective for information retrieval, their approach differs from ours in fundamental ways.
The focus of our approach is uncovering hierarchical relationships within topic structures across the dataset,
rather than enhancing document retrieval.
We do not discard high-support concepts,
as they provide crucial insights into overarching themes.
% Instead, we implicitly filter out low-support concepts through the TITANIC approach,
% ensuring that only representative structures are retained for analysis.

% FCA Topic Modeling
Our approach builds on the work of the authors of \cite{angelov_topic_modeling_2024} who propose extracting two relational structures from topic models.
The first structure is a weighted document-topic relation, i.e., which topics $t_i$ are present in which documents $d_j$.
The second structure is a weighted topic-word relation, i.e., which words $w_k$ are present in which topics $t_i$.
These relations are subsequently used to derive incidence matrices,
as visualized in Tables \ref{tab:doc-topic} and \ref{tab:term-topic}.
% Each row in Table \ref{tab:doc-topic} represents the topics assigned to a document,
% while each row in Table  \ref{tab:term-topic} captures the topics in which a term appears.

\begin{table}[t]
    \centering
    \begin{minipage}{0.5275\textwidth}
        \centering
        \caption{Document-topic relation}
        \label{tab:doc-topic}
        \begin{tabular}{|c|c|c|c|c|}
            \hline
            \backslashbox{Documents}{Topics} & $t_1$     & $t_2$     & $\cdots$ & $t_n$     \\ \hline
            $d_1$                            & $w_{1,1}$ & $w_{1,2}$ & $\cdots$ & $w_{1,n}$ \\ \hline
            $d_2$                            & $w_{2,1}$ & $w_{2,2}$ & $\cdots$ & $w_{2,n}$ \\ \hline
            $\vdots$                         & $\vdots$  & $\vdots$  & $\ddots$ & $\vdots$  \\ \hline
            $d_l$                            & $w_{l,1}$ & $w_{l,2}$ & $\cdots$ & $w_{l,n}$ \\ \hline
        \end{tabular}
    \end{minipage}%
    \hfill
    \begin{minipage}{0.465\textwidth}
        \centering
        \caption{Term-topic relation}
        \label{tab:term-topic}
        \begin{tabular}{|c|c|c|c|c|}
            \hline
            \backslashbox{Terms}{Topics} & $t_1$           & $t_2$           & $\cdots$ & $t_n$           \\ \hline
            $s_1$                        & $\hat{w}_{1,1}$ & $\hat{w}_{1,2}$ & $\cdots$ & $\hat{w}_{1,n}$ \\ \hline
            $s_2$                        & $\hat{w}_{2,1}$ & $\hat{w}_{2,2}$ & $\cdots$ & $\hat{w}_{2,n}$ \\ \hline
            $\vdots$                     & $\vdots$        & $\vdots$        & $\ddots$ & $\vdots$        \\ \hline
            $s_l$                        & $\hat{w}_{l,1}$ & $\hat{w}_{l,2}$ & $\cdots$ & $\hat{w}_{l,n}$ \\ \hline
        \end{tabular}
    \end{minipage}
\end{table}

% Since both Tables \ref{tab:doc-topic} and \ref{tab:term-topic} contain values in the range $[0,1]$,
% the authors propose additional processing steps to derive binary relations.
To compute a binary relation from the document-topic relation, they apply a thresholding approach,
while for the term-topic relation, they retain only the top $n$ words per topic.
Subsequently, the authors compute the formal concepts of high support.
Their approach enables a structured representation of topic hierarchies,
which they demonstrate on academic papers from machine learning conferences.
The approach lacks the possibility to compare subsets of the whole dataset, which might stem from an organization of files into directories.
% However, since the papers represent independent documents,
% their method lacks the ability to aggregate information across documents to derive insights into broader thematic structures.
% how our approach differs
By contrast, our approach addresses this limitation by explicitly incorporating cross-directory topic structures,
allowing for a more comprehensive analysis of thematic relationships spanning multiple data sources.
Furthermore, our method is not restricted to only textual data.

% Additionally, our methodology integrates text extracted from images, further enhancing the depth of topic modeling beyond purely textual data.
